% Part: proof-theory
% Chapter: proof-search
% Section: search-algorithm

\documentclass[../../../include/open-logic-section]{subfiles}

\begin{document}

\olfileid{pt}{ps}{sal}

\olsection{\Log{G3c}-র প্রমাণ-অন্বেষণ অ্যালগরিদম}

\Log{G3c}-র জন্য একটি প্রমাণ-অন্বেষণ অ্যালগরিদম বর্ণনা করি। এটিই একমাত্র
সম্ভব অ্যালগরিদম নয়, এমনকি সবচেয়ে দক্ষও নয়। তবে এটি সঠিক:
$\Gamma \Sequent \Delta$-র !!a{proof} থাকলে অ্যালগরিদমটি
একসময় থামবে এবং !!a{proof} খুঁজে পাবে। এটি ব্যর্থতার ক্ষেত্রেও
নির্ভরযোগ্য: !!a{proof} না পেয়ে থামলে, কিংবা কখনো না থামলে,
$\Gamma \Sequent \Delta$ শুরু থেকেই অবৈধ।

অ্যালগরিদমটির একটি অপরিহার্য বৈশিষ্ট্য হলো, $\Gamma \Sequent \Delta$-তে
বা প্রমাণ-অন্বেষণে তৈরি কোনো সিকোয়েন্টে থাকা প্রত্যেক !!{formula}
যতবার প্রয়োজন ততবার “হ্রাস” পায়; অর্থাৎ উল্টো দিকে প্রয়োগ করা
অনুমানবিধির প্রধান সূত্র হিসেবে নেওয়া হয়। এই বৈশিষ্ট্যের নাম
\emph{ন্যায্যতা}।

অ্যালগরিদমটি পর্যায়ক্রমে চলে। প্রতিটি পর্যায়ের ফল সিকোয়েন্টের একটি
বৃক্ষ। পরের পর্যায়ে, এখনও স্বতঃসিদ্ধ নয় এমন প্রত্যেক সর্বোচ্চ
সিকোয়েন্টে !!a{formula} হ্রাস করে নতুন বৃক্ষ তৈরি হয়।

$0$-তম পর্যায়ে $\Gamma \Sequent \Delta$ দিয়ে শুরু করি। ন্যায্যতা
নিশ্চিত করতে $\Gamma$ ও $\Delta$-র !!{formula}s-কে সূচক হিসেবে
এমন সংখ্যা দিই যাতে ভিন্ন !!{formula}s ভিন্ন সংখ্যা পায়। যেমন, বাঁ
থেকে ডান দিকে !!{formula}s গুনে সংখ্যা দেওয়া যায়। সংখ্যাগুলি ঊর্ধ্বলিপিতে
লিখি। যেমন, $!A, !B \Sequent !C$-র !!a{proof} খুঁজলে
$0$-তম পর্যায়ের ফল $!A^1, !B^2 \Sequent !C^3$। এমন
সূচকযুক্ত সিকোয়েন্টে ঊর্ধ্বলিপি হিসেবে থাকা সবচেয়ে বড়~$n$ হলো তার
\emph{সর্বোচ্চ সূচক} (উদাহরণে~$3$)।

সিকোয়েন্টে পরিমাণায়ক থাকলে, ডানে $\lexists[x][!A(x)]$ বা বাঁয়ে
$\lforall[x][!A(x)]$ হ্রাসে কোন পদ ব্যবহার করব, তাও জানতে হবে।
শুধু !!{constant}s $\Obj c_0$, $\Obj c_1$, $\Obj c_2$, \dots,
এবং~$\Gamma$ ও $\Delta$-তে থাকা ফাংশন-চিহ্ন দিয়ে গড়া পদই
প্রয়োজন। এই সব পদের একটি নির্দিষ্ট গণনাক্রম ধরে নিই, যাতে যেমন
“$\Gamma \Sequent \Delta$-তে নেই এমন প্রথম !!{constant}” বলা যায়।
অ্যালগরিদমে উৎপন্ন বৃক্ষের প্রতিটি সূচকযুক্ত সিকোয়েন্টের সঙ্গে
!!{constant}s-এর একটি সেট যুক্ত থাকে। $0$-তম পর্যায়ের সিকোয়েন্টে
যুক্ত !!{constant}s-এর সেট হলো সেখানে থাকা সব !!{constant}s-এর
সেট (যদি থাকে); কোনো !!{constant}s না থাকলে সেটি $\{\Obj c_0\}$।

ধরা যাক, $n$-তম পর্যায়ে অ্যালগরিদম !!{constant}s-এর সংশ্লিষ্ট
সেটসহ সূচকযুক্ত সিকোয়েন্টের একটি বৃক্ষ তৈরি করেছে। $n+1$-তম
পর্যায়ে প্রতিটি সর্বোচ্চ সিকোয়েন্টের জন্য নিচের কাজ করি।

কোনো পারমাণবিক !!a{formula}~$!A$ সিকোয়েন্টের বাঁ ও ডান উভয় দিকেই থাকলে,
অথবা বাঁ দিকে~$\lfalse$ থাকলে, কিছু করি না: এমন সিকোয়েন্টের
\Log{G3c}-তে !!{proof}s আছে, তাই এই সর্বোচ্চ সিকোয়েন্টে
আর কিছু করার নেই।

সিকোয়েন্টে অ-পারমাণবিক !!{formula} না থাকলে অ্যালগরিদম বন্ধ করি।
$\Gamma \Sequent \Delta$-র কোনো প্রমাণ নেই।

নয়তো সবচেয়ে ছোট সূচক~$i$-যুক্ত অ-পারমাণবিক !!{formula}~$!A$
বাছি। আলোচ্য সর্বোচ্চ সিকোয়েন্টটি তখন হয়
$!A^i, \Pi \Sequent \Lambda$, নয়
$\Pi \Sequent \Lambda, !A^i$। তার সর্বোচ্চ সূচক $k$ এবং
তার সঙ্গে যুক্ত !!{constant}s-এর সেট $C$ ধরি।

$!A^i$-কে “হ্রাস” করি: $!A$-র !!{main operator} এবং $!A^i$
বাঁয়ে না ডানে আছে, তা অনুযায়ী নির্ধারিত অনুমানবিধি দিয়ে নতুন
সর্বোচ্চ সিকোয়েন্ট তৈরি করি।

\begin{enumerate}
  \item $!A \ident !B \land !C$ এবং বাঁ দিকে আছে:
  $(!B \land !C)^i, \Pi \Sequent \Lambda$-র উপরে একটি নতুন
  সর্বোচ্চ সিকোয়েন্ট যোগ করি:
  \[
  \Axiom$!B^{k+1}, !C^{k+2}, \Pi \fCenter \Lambda$
  \RightLabel{\LeftR{\land}}
  \UnaryInf$(!B \land !C)^i, \Pi \fCenter \Lambda$
  \DisplayProof
  \]
  নতুন সিকোয়েন্টের সঙ্গে যুক্ত !!{constant}s-এর সেট~$C$।
  \item $!A \ident !B \land !C$ এবং ডান দিকে আছে:
  $\Pi \Sequent \Lambda, (!B \land !C)^i$-র উপরে দুটি
  নতুন সর্বোচ্চ সিকোয়েন্ট যোগ করি:
  \[
  \Axiom$\Pi \fCenter \Lambda, !B^{k+1}$
  \Axiom$\Pi \fCenter \Lambda, !C^{k+1}$
  \RightLabel{\RightR{\land}}
  \BinaryInf$\Pi \fCenter \Lambda, (!B \land !C)^i$
  \DisplayProof
  \]
  নতুন সিকোয়েন্টগুলির সঙ্গে যুক্ত !!{constant}s-এর সেট~$C$।

  \item $!A \ident \lnot !B$, $!A \ident !B \lor !C$,
  এবং $!A \ident !B \lif !C$ ক্ষেত্রগুলি একইরকম;
  অনুশীলনের জন্য রাখা হলো।

  \item $!A \ident \lexists[x][!B(x)]$ এবং বাঁ দিকে আছে:
  $\lexists[x][!B(x)]^i, \Pi \Sequent \Lambda$-র উপরে একটি
  নতুন সর্বোচ্চ সিকোয়েন্ট যোগ করি:
  \[
  \Axiom$!B(c)^{k+1}, \Pi \fCenter \Lambda$
  \RightLabel{\LeftR{\lexists}}
  \UnaryInf$\lexists[x][!B(x)]^i, \Pi \fCenter \Lambda$
  \DisplayProof
  \]
  এখানে $c$ হলো $C$-তে নেই এমন প্রথম !!{constant}।
  নতুন সিকোয়েন্টের সঙ্গে যুক্ত !!{constant}s-এর সেট
  $C \cup \{c\}$।
  \item $!A \ident \lexists[x][!B(x)]$ এবং ডান দিকে আছে:
  $\Pi \Sequent \Lambda, \lexists[x][!B(x)]^i$-র উপরে একটি
  নতুন সর্বোচ্চ সিকোয়েন্ট যোগ করি:
  \[
  \Axiom$\Pi \fCenter \Lambda, \lexists[x][!B(x)]^{k+1}, !B(t)^{k+2}$
  \RightLabel{\RightR{\lexists}}
  \UnaryInf$\Pi \fCenter \Lambda, \lexists[x][!B(x)]^i$
  \DisplayProof
  \]
  এখানে $t$ হলো $C$-র !!{constant}s দিয়েই গড়া এবং
  এই সিকোয়েন্টের নিচে ডানে $\lexists[x][!B(x)]$ হ্রাসে
  এখনও ব্যবহৃত হয়নি এমন প্রথম পদ (এমন সব পদ ব্যবহৃত
  হলে ওই সেটের প্রথম ধ্রুবক নিই)। নতুন সিকোয়েন্টের সঙ্গে যুক্ত
  !!{constant}s-এর সেট $C$।
  % BN-SRC-601 TARGET-CORRECTION-BEGIN
  \par\smallskip\noindent\textnormal{\small\textbf{উৎস-সংশোধন (BN-SRC-601):} ডান অস্তিত্ব-বিধির সিদ্ধান্তে বাদ-পড়া সূচক ফিরেছে; বাঁ বিধির ভুল লেবেলের বদলে ডান বিধির লেবেল বসেছে।}
  % BN-SRC-601 TARGET-CORRECTION-END
  \item $!A \ident \lforall[x][!B(x)]$ ক্ষেত্রগুলি একইরকম;
  অনুশীলনের জন্য রাখা হলো।
\end{enumerate}
% BN-SRC-600 TARGET-CORRECTION-BEGIN
\par\smallskip\noindent\textnormal{\small\textbf{উৎস-সংশোধন (BN-SRC-600):} $k$ বিদ্যমান সর্বোচ্চ সূচক; উৎসের নতুন সূত্রগুলিতে $k$ আবার ব্যবহৃত হওয়ায় ন্যায্যতার যুক্তি ব্যর্থ হতে পারে। চারটি বিধি-ক্ষেত্রে সব নতুন সূচক $k$-র পরে রাখা হয়েছে; পৃথক শাখায় একই নতুন সূচক গ্রহণযোগ্য।}
% BN-SRC-600 TARGET-CORRECTION-END
% BN-SRC-888 TARGET-CORRECTION-BEGIN
\par\smallskip\noindent\textnormal{\small\textbf{পাঠ-সংশোধন (BN-SRC-888):} উপলব্ধ পদগুলি ফুরোলে ইতিমধ্যে নথিভুক্ত ধ্রুবকই আবার নেওয়া হয়। উৎসের নির্দিষ্ট ধ্রুবকটি ওই সেটে না-ও থাকতে পারে; সেটি না লিখে ব্যবহার করলে পরের নতুন-ধ্রুবক বিধির শর্ত নষ্ট হয়।}
% BN-SRC-888 TARGET-CORRECTION-END

$n+1$-তম পর্যায়ে কোনো সর্বোচ্চ সিকোয়েন্টের উপরে নতুন সিকোয়েন্ট
যোগ না হলে অ্যালগরিদম সফলভাবে শেষ হয়। তখন $n+1$-তম
পর্যায়ের বৃক্ষের সব সূচক মুছে $\Gamma \Sequent \Delta$-র
!!a{proof} পাই। সব সর্বোচ্চ সিকোয়েন্ট তখন
$!A, \Pi \Sequent \Lambda, !A$ বা
$\lfalse, \Pi \Sequent \Lambda$ আকারের। প্রথম আকারের অভিন্ন সূত্রটি পারমাণবিক। সব অনুমানই
\Log{G3c}-র সঠিক অনুমান। প্রস্তাবনামূলক অনুমানের ক্ষেত্রে
এটি স্পষ্ট। লক্ষ করি, প্রতি ধাপে একটি সিকোয়েন্টের সঙ্গে যুক্ত
!!{constant}s-এর সেট~$C$ সেই সিকোয়েন্টে থাকা সব
!!{constant}s ধারণ করে। তাই \LeftR{\lexists} ও
\RightR{\lforall} ব্যবহার করে হ্রাসের সময় আইগেনচল-শর্ত
পূরণ হয়: আইগেনচল~$c$ ছিল !!a{constant}, যা
সিদ্ধান্তের সঙ্গে যুক্ত সেট~$C$-তে ছিল না; তাই $c$
সিদ্ধান্তেও ছিল না। সুতরাং:

% BN-SRC-920 TARGET-CORRECTION-BEGIN
\par\smallskip\noindent\textnormal{\small\textbf{উৎস-সংশোধন (BN-SRC-920):} G3c-র অভেদ-স্বতঃসিদ্ধে অভিন্ন সূত্রটি পারমাণবিক। যৌগিক সূত্র উভয় দিকে থাকলেই একটি পংক্তিকে স্বতঃসিদ্ধ ধরা যায় না; সফল অন্বেষণের শেষ পংক্তির বর্ণনাতেও ওই শর্ত রাখা হয়েছে।}
% BN-SRC-920 TARGET-CORRECTION-END
\begin{prop}
\Log{G3c}-র প্রমাণ-অন্বেষণ অ্যালগরিদম সঠিক: এটি সফলভাবে শেষ
হলে প্রাথমিক সিকোয়েন্ট~$\Gamma \Sequent \Delta$-র !!a{proof}
থাকে।
\end{prop}





\end{document}
